% 逻辑陪集可区分性(LCD)理论笔记 —— 主文件
% 编译:  cd theory && latexmk -xelatex main.tex     (需要 ctex;见 theory/README.md)
\documentclass[UTF8,a4paper,11pt]{ctexart}

\usepackage{amsmath,amssymb,amsthm,mathtools}
\usepackage{booktabs}
\usepackage{longtable}
\usepackage[margin=2.6cm]{geometry}
\usepackage[colorlinks=true,linkcolor=blue!50!black,citecolor=blue!50!black,urlcolor=blue!50!black]{hyperref}
\usepackage{xcolor}
\usepackage{enumitem}

% ---------- 定理环境 ----------
\theoremstyle{plain}
\newtheorem{theorem}{定理}[section]
\newtheorem{lemma}[theorem]{引理}
\newtheorem{proposition}[theorem]{命题}
\newtheorem{corollary}[theorem]{推论}
\newtheorem{conjecture}[theorem]{猜想}
\theoremstyle{definition}
\newtheorem{definition}[theorem]{定义}
\newtheorem{problem}[theorem]{开放问题}
\theoremstyle{remark}
\newtheorem{remark}[theorem]{注}
\newtheorem{observation}[theorem]{数值观察}

% ---------- 状态标签(可 grep):PROVED / SKETCH / NUMERICAL / CONJECTURE / PLAN ----------
%  PROVED     笔记中给出完整论证(仍需作者逐行复核,见 theory/README.md)
%  SKETCH     论证框架完整,个别步骤引用标准结果或尚未展开
%  NUMERICAL  由 theory/checks 中的精确计算观察到,尚无证明
%  CONJECTURE 明确的猜想
%  PLAN       只有陈述与证明路线
\newcommand{\st}[1]{{\small\textsc{[\textcolor{red!60!black}{#1}]}}}

% ---------- 记号 ----------
\newcommand{\F}{\mathbb{F}_2}
\newcommand{\Prob}{\mathbb{P}}
\newcommand{\E}{\mathbb{E}}
\newcommand{\err}{\varepsilon^{\star}}           % Bayes (最大似然) 错误率 = 逻辑失败概率
\newcommand{\pL}{p_{\mathrm{L}}}
\newcommand{\Dep}{\mathrm{Dep}}
\newcommand{\Nn}{\mathcal{N}}
\newcommand{\supp}{\operatorname{supp}}
\newcommand{\wt}{\operatorname{wt}}
\newcommand{\BC}{\mathfrak{B}}                    % Bhattacharyya / Hellinger 系数
\renewcommand{\Pr}{\Prob}

\title{逻辑陪集可区分性的理论笔记\\\large 单调性引理、局域汇率上界、预序、条件认证与实验设计}
\author{LCD 项目\ (theory 线)}
\date{\today\ \ ——\ 草稿,持续更新}

\begin{document}
\maketitle

\begin{abstract}\noindent
本笔记是 \texttt{ABSTRACT.md} 中四类理论结果的工作底稿。每个命题都带有状态标签
(\st{PROVED}、\st{SKETCH}、\st{NUMERICAL}、\st{CONJECTURE}、\st{PLAN}),并在 §\ref{sec:map}
中与仓库里的数值检验一一对应:\textbf{理论带路,数值检验理论}。
已写出完整证明的结果:(i) 三类自由操作下的单调性;(ii) 擦除/泄漏/突发事件在模型中的表示及适用边界;
(iii) 局域汇率的精确上界
$R_{e\to p}\le (3/4-p_0)/(1-e_0)$;(iv) 码容量泡利加擦除类中局域自由操作预序的完整刻画;
(v) 代数可计算的有限样本认证(码容量)。其余为带证明路线的陈述。
\end{abstract}

\tableofcontents

\section{框架:线性解码结构、可模拟约化与单调性}\label{sec:framework}

\subsection{线性解码结构与噪声实例}

为了让同一套证明同时覆盖码容量噪声和电路级噪声(检测器误差模型),我们只用到\emph{线性性}。

\begin{definition}[线性解码结构]\label{def:structure}
一个线性解码结构是三元组 $\mathfrak D=(V,\sigma,L)$:$V$ 是有限维 $\F$-线性空间(\emph{故障空间}),
$\sigma:V\to\F^{r}$ 是\emph{检测器映射}(综合征),$L:V\to\F^{q}$ 是\emph{可观测量映射},均为线性映射。
\end{definition}

\begin{definition}[稳定子码的结构]\label{def:code-structure}
设 $S$ 为 $n$ 比特稳定子码的稳定子群($n-k$ 个独立生成元 $g_i$),泡利算符模相位记为 $V=\F^{2n}$,
辛内积 $\langle u,v\rangle=u_x\!\cdot v_z+u_z\!\cdot v_x$。取
\[
\sigma(E)=\big(\langle E,g_i\rangle\big)_{i=1}^{n-k},\qquad
L(E)=\big(\langle E,\bar X_j\rangle,\langle E,\bar Z_j\rangle\big)_{j=1}^{k},
\]
其中 $\bar X_j,\bar Z_j$ 是一组逻辑算符(辛基)。映射 $E\mapsto(\sigma(E),L(E))$ 满射,其核恰为 $S$。
电路级噪声:$V=\F^{\#\text{故障位置}}$,$\sigma$ 为检测器,$L$ 为可观测量翻转,线性性来自
Clifford 电路中泡利故障的线性传播。
\end{definition}

\begin{definition}[噪声实例与陪集区分问题]\label{def:instance}
噪声实例是随机对 $(E,F)$:$E\in V$ 是故障向量,$F\in\mathfrak F$($\mathfrak F$ 有限)是经典\emph{边信息}
(擦除标记、泄漏检测标记、事件标记等),联合分布 $P_{EF}$ 为解码器所知。解码器的观测为
$Y=(\sigma(E),F)$,任务是猜测 $L(E)$。解码器是(可随机化的)映射 $\hat g(Y)$。定义
\[
\err(P):=\min_{\hat g}\Prob[\hat g(Y)\neq L(E)]
=1-\sum_{y}\max_{\ell}\Prob[L=\ell,\,Y=y].
\]
最小值由最大似然(陪集)解码达到,$\err(P)=\pL$ 即最大似然下的逻辑失败概率。
\end{definition}

\begin{lemma}[解码即标签猜测]\label{lem:decoding}
对稳定子码(定义~\ref{def:code-structure}),设解码器看到 $y=(s,f)$ 后施加恢复算符 $R$,$\sigma(R)=s$。
则恢复成功(即 $E+R\in S$,逻辑恒等)当且仅当 $L(R)=L(E)$;且对每个 $(s,\ell)$ 存在 $R$ 使
$\sigma(R)=s,L(R)=\ell$。因此恢复失败概率的最小值就是 $\err(P)$。\hfill\st{PROVED}
\end{lemma}
\begin{proof}
$\sigma(E+R)=0$,故 $E+R$ 属于归一化子,且 $E+R\in S\iff L(E+R)=0\iff L(R)=L(E)$
(因 $(\sigma,L)$ 的联合核为 $S$)。$(\sigma,L)$ 满射给出第二句。
\end{proof}

\begin{remark}[标签选取无关]
把 $L$ 换成 $L'(E)=\pi_{\sigma(E)}(L(E))$,其中 $\pi_s$ 是依赖综合征的标签置换(例如换一个参考提升 $\rho(s)$),
$\err$ 不变:对固定 $y$ 取最大的 $\ell$ 的值不受置换影响。数值上我们直接使用原始标签
$(\langle E,\bar X\rangle,\langle E,\bar Z\rangle)$(见 \texttt{theory/checks})。
\end{remark}

对一族噪声模型 $\Nn=(P_d)_{d}$(码距 $d$ 为指标),记
$\alpha_-(\Nn)=\liminf_{d\to\infty}-\tfrac1d\ln\err(P_d)$,$\alpha_+(\Nn)=\limsup_{d\to\infty}-\tfrac1d\ln\err(P_d)$;
二者相等时记为压制指数 $\alpha$。

\subsection{可模拟约化}

\begin{definition}[可模拟约化]\label{def:reduction}
设 $(L,Y)$ 与 $(L',Y')$ 是取值于有限集的随机对。一个\emph{可模拟约化}
$(L,Y)\rightsquigarrow(L',Y')$ 是三元组 $(\kappa,g,h)$:
$\kappa(\cdot\mid y)$ 是从 $Y$ 到\emph{处理器随机性} $\xi$ 的马尔可夫核;$g(y,\xi)$ 与 $h(\ell,y,\xi)$ 为函数,满足

\smallskip
(i) 对每个 $(y,\xi)$,$\ell\mapsto h(\ell,y,\xi)$ 是标签集上的双射;

(ii) 令 $\xi\sim\kappa(\cdot\mid Y)$(在给定 $Y$ 时与 $L$ 条件独立),则
$\big(h(L,Y,\xi),\,g(Y,\xi)\big)$ 与 $(L',Y')$ 同分布。
\end{definition}

\begin{theorem}[单调性]\label{thm:monotone}
若 $(L,Y)\rightsquigarrow(L',Y')$,则 $\err(L\mid Y)\le\err(L'\mid Y')$。\hfill\st{PROVED}
\end{theorem}
\begin{proof}
设 $\hat g'$ 是 $(L',Y')$ 上的任一解码器。构造 $(L,Y)$ 上的随机化解码器 $\hat g$:
观察 $y$,抽 $\xi\sim\kappa(\cdot\mid y)$,输出 $h^{-1}(\hat g'(g(y,\xi));y,\xi)$,其中 $h^{-1}$ 对标签变量取逆(由 (i) 存在)。
则
\[
\{\hat g(Y)=L\}=\{\hat g'(g(Y,\xi))=h(L,Y,\xi)\}=\{\hat g'(Y'_{\rm sim})=L'_{\rm sim}\},
\]
其中 $(L'_{\rm sim},Y'_{\rm sim})=(h(L,Y,\xi),g(Y,\xi))$ 由 (ii) 与 $(L',Y')$ 同分布。故
$\Prob[\hat g(Y)\neq L]=\Prob[\hat g'(Y')\neq L']$。随机化解码器的错误率是确定性解码器错误率的平均,
所以存在确定性解码器的错误率不超过它。对 $\hat g'$ 取最优即得结论。
\end{proof}

\begin{remark}[数据处理不等式的形式]
$1-\err(L\mid Y)=2^{-H_{\min}(L\mid Y)}$($H_{\min}$ 为条件最小熵)。当 $h=\mathrm{id}$
(只处理观测)时,定理~\ref{thm:monotone} 就是 $H_{\min}(L\mid Y)\le H_{\min}(L\mid g(Y,\xi))$,
即条件最小熵(Sibson--R\'enyi $\infty$ 阶互信息)的数据处理不等式。其余 R\'enyi/Hellinger 型泛函同样单调,
见命题~\ref{prop:hellinger}。
\end{remark}

\subsection{三类自由操作}

下面三类操作都是定义~\ref{def:reduction} 的特例,从而由定理~\ref{thm:monotone} 得到单调性。
设 $(E,F)$ 是线性解码结构 $\mathfrak D$ 上的噪声实例。

\begin{description}[leftmargin=2.2em,style=nextline]
\item[(F1) 综合征粗粒化]
$Y'=g(Y,\xi)$,$h=\mathrm{id}$,$g$ 为任意(可随机)函数,例如合并检测器、取子集。\\
\emph{特例}:丢弃标记(\textbf{(F1$'$)}):$F'=\varphi(F,\xi)$,仍是噪声实例 $(E,F')$。
\item[(F2) 叠加可经典采样的独立噪声]
设 $(E',F'')$ 是与 $(E,F)$ 独立、分布已知的随机对。新实例为
$(E+E',\,(F,F''))$(或经 (F1$'$) 进一步粗粒化的版本)。
取 $\xi=(E',F'')$,$g(y,\xi)=(\sigma(E')+s,\,f,\,F'')$,$h(\ell,y,\xi)=\ell+L(E')$。
\item[(F2$'$) 依赖已观测标记的独立噪声]
$\xi\sim\kappa(\cdot\mid y)$ 可依赖观测 $y=(s,f)$;例如只对未被标记的比特叠加去极化噪声。
$h(\ell,y,\xi)=\ell+L(N)$,$N$ 为处理器自己抽出的叠加故障。
\end{description}

\begin{proposition}[(F1)(F2)(F2$'$) 为可模拟约化]\label{prop:free-ops}
设 $(E,F)$ 为噪声实例,按 (F2) 或 (F2$'$) 得到 $(E_{\rm new},F_{\rm new})$,其中叠加的故障 $N$ 满足:
给定 $F$ 时 $N$ 与 $E$ 条件独立(特别地,(F2) 中 $N=E'$ 与 $(E,F)$ 独立)。则
\[
\err(E,F)\;\le\;\err(E+N,\,F_{\rm new}),
\]
其中 $F_{\rm new}$ 是 $(F,\text{附加标记})$ 的任一(随机)函数。特别地,叠加噪声永不降低 $\err$。
\hfill\st{PROVED}
\end{proposition}
\begin{proof}
新问题的标签为 $L(E+N)=L(E)+L(N)$(线性性),观测为 $(\sigma(E)+\sigma(N),F_{\rm new})$。
处理器由 $(\sigma(E),F)$ 与自己抽取的 $(N,\text{附加标记})$ 算出新观测,并用已知的 $L(N)$ 做标签平移。
这是定义~\ref{def:reduction} 的约化,其中 (i) 因平移是双射,(ii) 因 $N$ 条件独立于 $E$ 且分布已知
给出正确联合分布。再应用定理~\ref{thm:monotone}。
\end{proof}

\begin{corollary}[指数的单调性]\label{cor:alpha}
若对每个 $d$,$P'_d$ 由 $P_d$ 经(同类型的)自由操作得到,则
$\alpha_\pm(\Nn')\le\alpha_\pm(\Nn)$。\hfill\st{PROVED}
\end{corollary}

\begin{remark}[独立性不可去掉]\label{rem:independence}
取 $N=E$(叠加的噪声是已有故障的拷贝):$E+N=0$,$\err=0<\err(E,F)$。所以
处理器必须能在\emph{不知道隐变量 $E$} 的情况下抽出 $N$——这正是“叠加\emph{独立}噪声”的含义。
第~\ref{sec:repr} 节据此划定模型类的边界。
\end{remark}

\begin{remark}[只对最优解码器成立]\label{rem:ml-only}
定理~\ref{thm:monotone} 刻画的是 Bayes 最优(最大似然)解码。对次优解码器(如固定权重的 MWPM),
叠加噪声并不保证提高其失败率的单调性,也没有现成的数据处理不等式;因此单调性的\emph{数值}检验必须使用精确 ML
(见 \texttt{theory/checks/exact\_small\_codes.py}),不能用 MWPM 结果代替。
反过来,由最优性始终有 $\err\le p_{\rm L}^{\rm MWPM}$,即 $\alpha_{\rm ML}\ge\alpha_{\rm MWPM}$。
\end{remark}

\section{擦除、泄漏与事件类噪声在模型中的表示}\label{sec:repr}

单调性(定理~\ref{thm:monotone})只对满足“可经典采样、与隐变量独立”的叠加成立
(注~\ref{rem:independence})。本节逐一说明各类物理噪声如何写成噪声实例 $(E,F)$,
以及单调性论证在什么假设下仍然成立、在哪里会失效。

\subsection{擦除}

\begin{proposition}[擦除 = 带标记的完全去极化]\label{prop:erasure}
设稳定子码,理想(无噪声)综合征提取。比特 $j$ 以概率 $e$ 被擦除并被标记($F_j=1$),擦除后该比特处于最大混态。
则在 $(E,F)$ 表示下,擦除等价于:$F_j=1$ 且 $E_j$ 在 $\{I,X,Y,Z\}$ 上均匀分布,与其他比特独立。
\hfill\st{PROVED}
\end{proposition}
\begin{proof}
替换为最大混态的信道 $\Lambda(\rho)=\operatorname{tr}(\rho)\,\mathbb 1/2=\tfrac14\sum_{P\in\{I,X,Y,Z\}}P\rho P$
(泡利扭曲恒等式)。于是整体信道是泡利酉算符的凸组合;综合征测量是泡利框架上的线性函数,
故整个过程等价于“随机泡利 + 经典标记”。
\end{proof}

\begin{remark}[保守性]
若实际擦除把比特重置到固定已知态(如 $\lvert 0\rangle$)或留下泄漏态,实验者可以额外对该位置施加一个随机泡利
(带标记)——这又是 (F2$'$) 型的自由操作,只会使问题变难。因此“最大混态模型”给出的
$\err$ 是\emph{真实噪声 $\err$ 的上界}(对真实噪声保守的一侧),而不是下界。
\end{remark}

\begin{corollary}[丢弃标记 = 去极化]\label{cor:discard}
对被擦除的比特丢弃标记 (F1$'$),得到未标记的均匀泡利错误,即速率 $3/4$ 的去极化(非平凡泡利的概率为 $3/4$)。
\end{corollary}

\subsection{码容量类 $\Nn_{\rm cc}(p,e)$ 及其上的三种操作}

\begin{definition}\label{def:Ncc}
$\Nn_{\rm cc}(p,e)$:每个数据比特独立,以概率 $e$ 被擦除(标记 $F_j=1$,$E_j$ 均匀),
否则以概率 $p$ 发生去极化($E_j=X,Y,Z$ 各 $p/3$)。记 $\lambda(p)=1-\tfrac43p$(去极化信道的泡利本征值),
$\mu(p,e)=(1-e)\lambda(p)$。注意 $\Dep_p$ 对 $p$ 是线性的,$\Dep_{3/4}$ 是均匀分布。
\end{definition}

\begin{lemma}[类内的基本自由操作]\label{lem:basic-ops}
下列操作把 $\Nn_{\rm cc}(p,e)$ 映到 $\Nn_{\rm cc}(p',e')$,且都是 (F1$'$)/(F2)/(F2$'$) 的特例:
\begin{enumerate}[label=(\alph*),leftmargin=2.2em]
\item \emph{叠加去极化}(F2$'$,只作用在未标记比特上):$p'=p\oplus q$,即 $\lambda(p')=\lambda(p)\lambda(q)$;$e'=e$。
\item \emph{叠加擦除}(F2):$p'=p$,$e'=1-(1-e)(1-\delta)$。
\item \emph{随机丢弃标记}(F1$'$):每个被标记比特独立地以概率 $\kappa\in[0,1]$ 丢弃标记。则
\[
e'=e(1-\kappa),\qquad p'=\frac{(1-e)p+\tfrac34 e\kappa}{1-e+e\kappa},\qquad
\lambda(p')=\frac{(1-e)\lambda(p)}{1-e+e\kappa}.
\]
\end{enumerate}
\hfill\st{PROVED}
\end{lemma}
\begin{proof}
(a)(b) 由去极化信道的复合与“擦除覆盖已有故障”的事实(均匀噪声与任何独立噪声之和仍均匀)直接得到;
(b) 中被新擦除的比特其 $E$ 变为均匀,标记并入。

(c) 每个比特独立处理。丢弃标记后:仍被标记者概率 $e(1-\kappa)$,此时 $E$ 均匀;未标记者以概率
$(1-e)$ 来自未擦除(分布 $\Dep_p$),以概率 $e\kappa$ 来自被丢弃标记的擦除(分布 $\Dep_{3/4}$)。
给定“未标记”,$E$ 的分布是 $\Dep_p$ 与 $\Dep_{3/4}$ 按权重 $(1-e):e\kappa$ 的混合;
因 $\Dep_p$ 对 $p$ 线性,混合仍是 $\Dep_{p'}$,$p'$ 即上式;比特之间独立,故结果仍在 $\Nn_{\rm cc}$ 内。
$\lambda$ 的公式由 $\lambda=1-\tfrac43p$ 的线性性得到。
\end{proof}

\subsection{泄漏}

物理上,泄漏比特离开计算子空间,在其寿命内使与之相互作用的门产生随机泡利错误并使测量结果随机。
我们把泄漏纳入模型类的\emph{假设}如下。

\begin{definition}[假设 (L):泄漏可经典采样]\label{def:L}
泄漏过程由随机\emph{泄漏事件集合} $\Lambda$(位置、起始时间、持续时间 $\le\tau$)描述,满足
\begin{enumerate}[label=(L\arabic*),leftmargin=2.8em]
\item 经计算子空间泡利扭曲后,泄漏对检测器的全部影响等价于故障向量 $E_\Lambda\in V$,其条件分布
$P(E_\Lambda\mid\Lambda)$ 已知(典型:被泄漏比特参与的每个门以概率 $1/2$ 施加均匀泡利);
\item $\Lambda$ 的分布已知,且与背景故障 $E_{\rm bg}$ \emph{独立};
\item 边信息 $F_\Lambda$ 是 $\Lambda$ 的已知(可能是部分)函数:全检测对应带标记,无检测对应 $F_\Lambda=\text{常数}$。
\end{enumerate}
\end{definition}

\begin{proposition}[泄漏下的单调性]\label{prop:leak}
在假设 (L) 下,总噪声实例为 $(E_{\rm bg}+E_\Lambda,\,(F_{\rm bg},F_\Lambda))$,而 $(E_\Lambda,F_\Lambda)$ 是与背景
独立、可经典采样的随机对。于是命题~\ref{prop:free-ops} 适用:\\
(1) 加入泄漏不降低 $\err$;\\
(2) 泄漏事件率 $r_\ell$ 的单调性:速率为 $r_\ell+r'_\ell$ 的泊松泄漏过程 = 速率 $r_\ell$ 的过程与独立的速率 $r'_\ell$ 过程之和
(泊松叠加定理;故障按 $\F$ 相加),故 $\err$ 对 $r_\ell$ 不降;\\
(3) 丢弃泄漏检测标记是 (F1$'$),使 $\err$ 不降。\hfill\st{PROVED}
\end{proposition}

\begin{remark}[假设 (L) 之外:论证失效的位置]\label{rem:L-fail}
(L2) 失败时,叠加层依赖隐变量,处理器无法独立采样(注~\ref{rem:independence}),单调性论证不成立。典型情形:
(i) 泄漏概率依赖比特所处的计算基态(例如 $\lvert1\rangle$ 更易泄漏),使 $\Lambda$ 与背景故障相关;
(ii) 泄漏比特对邻居的作用是相干旋转而非随机泡利,(L1) 的扭曲近似不成立;
(iii) 泄漏与测量/重置错误通过同一物理机制耦合。这些都被列为\emph{模型类假设},不在本笔记的证明范围内,
并应与 ABSTRACT 的“范围与局限”一致:相干误差不在模型类内。
(L1) 中“扭曲后为泡利”的近似本身是可由实验检验的假设,应与相干误差一并列为模型类假设。
\end{remark}

\subsection{局域突发事件与芯片尺度事件}

\begin{definition}[加性事件层]\label{def:events}
设背景 $(E_{\rm bg},F_{\rm bg})$ 如上。\emph{局域突发层}:时空中的泊松点过程,强度 $r\,\mu(dz\,dt)$,
每个事件(中心 $z$、起始 $t$)在半径 $R$、持续 $T_b$ 的时空区域内的每个故障位置独立地施加速率 $p_b$ 的去极化噪声;
\emph{芯片尺度层}:速率 $r_c$ 的泊松事件,每个事件在持续 $\tau_c$ 内对全芯片位置施加某个已知故障律 $K_c$。
各层与背景相互独立,故障按 $\F$ 相加(\emph{加性语义},与 README P2 中“检测事件按位异或叠加”一致)。
\end{definition}

\begin{proposition}[速率方向的单调性]\label{prop:rate-mono}
在定义~\ref{def:events} 的模型类中,$\err$ 对 $p$、$e$、突发率 $r$、芯片事件率 $r_c$ 均不降。
\hfill\st{PROVED}
\end{proposition}
\begin{proof}
强度 $r+r'$ 的泊松过程等于强度 $r$ 与强度 $r'$ 的独立过程之并(叠加定理),且加性语义下各事件的故障按 $\F$ 相加,
因此增大 $r$ 恰是叠加一个独立的、可经典采样的突发层;$r_c$、$p$、$e$ 同理(引理~\ref{lem:basic-ops})。
应用命题~\ref{prop:free-ops}。
\end{proof}

\begin{remark}[形状参数 $(R,T_b,p_b,\tau_c)$ 的单调性\emph{不}由自由操作给出]\label{rem:shape}
增大 $p_b$(或 $R,T_b$)相当于在\emph{已发生的突发区域内}再叠加噪声。突发事件对解码器不可见(无标记)时,
该附加层依赖隐变量(事件位置),处理器无法独立采样,论证失败。
若突发事件\emph{带标记}(事件的位置与时间进入 $F$),则附加层依赖观测到的 $F$,属 (F2$'$),单调性成立。
无标记情形下关于形状参数的单调性是开放问题(\ref{prob:O1})。这直接决定认证定理
(§\ref{sec:cert})能否对形状不确定性给出稳健上界。
\end{remark}

\section{The local exchange-rate bound and the preorder on the Pauli-plus-erasure class}\label{sec:exchange}

The main line of this section: \emph{discarding an erasure flag is a free operation} (Corollary~\ref{cor:discard}), which gives a rigorous upper bound on the exchange rate (Theorem~\ref{thm:rate}),
and in the simplest subclass $\Nn_{\rm cc}(p,e)$ we completely characterize the preorder induced by local free operations (Theorem~\ref{thm:preorder}).
The gap between the true exchange rate and the bound is one of the core research questions of the project (\S\ref{sec:gap}).

\subsection{Partial flag discarding and the exchange-rate bound}

\begin{theorem}[Partial flag-discarding inequality]\label{thm:partial}
Let $0\le p\le\tfrac34$, $0\le e_0<1$, $0\le\delta\le1-e_0$. For every code distance $d$,
\[
\err_d\big(\Nn_{\rm cc}(p,\,e_0+\delta)\big)\;\le\;\err_d\big(\Nn_{\rm cc}(p'',\,e_0)\big),\qquad
p''=p+\delta\,\frac{\tfrac34-p}{1-e_0}.
\]
Consequently $\alpha_\pm(p,e_0+\delta)\ge\alpha_\pm(p'',e_0)$.\hfill\st{PROVED}
\end{theorem}
\begin{proof}
Apply Lemma~\ref{lem:basic-ops}(c) with $e=e_0+\delta$ and $\kappa=\delta/(e_0+\delta)$ (trivial if $e_0+\delta=0$).
Then $e'=e(1-\kappa)=e_0$ and $e\kappa=\delta$, so
\[
p'=\frac{(1-e_0-\delta)p+\tfrac34\delta}{1-e_0}=p+\delta\frac{\tfrac34-p}{1-e_0}=p''.
\]
This operation is (F1$'$), so Theorem~\ref{thm:monotone} gives $\err(p,e_0+\delta)\le\err(p'',e_0)$;
the exponent version follows from Corollary~\ref{cor:alpha}.
\end{proof}

\begin{theorem}[Exchange-rate bound]\label{thm:rate}
Let $c(p_0,e_0)=\dfrac{3/4-p_0}{1-e_0}$.

(a) If the right derivatives of $\err_d$ with respect to $p$ and $e$ exist at $(p_0,e_0)$, then
\[\partial_e\err_d\le c\,\partial_p\err_d,\]
and when $\partial_p\err_d>0$, $R^{(d)}:=\partial_e\err_d/\partial_p\err_d\le c(p_0,e_0)$.

(b) Suppose the limit $\alpha=\alpha_-=\alpha_+$ exists, has right derivatives with respect to $p$ and $e$ at $(p_0,e_0)$, and $\partial_p\alpha<0$. Then the local exchange rate satisfies
\[
R_{e\to p}(p_0,e_0)=-\frac{\partial_e\alpha}{\partial_p\alpha}\;\le\;\frac{3/4-p_0}{1-e_0}.
\]
\hfill\st{PROVED}
\end{theorem}
\begin{proof}
In Theorem~\ref{thm:partial}, $p''=p_0+c\,\delta$ is linear in $\delta$. (a): subtract $\err_d(p_0,e_0)$ from both sides and divide by $\delta\downarrow0$;
the left side tends to $\partial_e\err_d$ and the right side to $c\,\partial_p\err_d$. (b): the same for $\alpha$ gives
$\partial_e\alpha\ge c\,\partial_p\alpha$; since $-\partial_p\alpha>0$, multiplying by $-1$ gives
$-\partial_e\alpha\le c(-\partial_p\alpha)$, and dividing by $-\partial_p\alpha$ gives the claim.
\end{proof}

\paragraph{Constitutional reading.}
$R_{e\to p}$ here is the \emph{marginal rate of a monotone} (\S\ref{sec:rates}); by Proposition~\ref{prop:marginal} every monotone of the local free order has marginal rate in $[0,c]$, and $c$ is attained by the free-order monotone $1-\mu$, so Theorem~\ref{thm:rate} cannot be improved using the free operations alone.

\begin{remark}[Relation to the heuristic ``$R\le3/4$'']\label{rem:34}
``Adding $\delta e$ of erasure $\to$ adding $\tfrac34\delta e$ of depolarizing noise'' is the value of $c$ at $(p_0,e_0)=(0,0)$. The exact $c$ has two corrections:
(i) the numerator $\tfrac34-p_0$: a newly erased qubit already carried depolarizing noise of probability $p_0$, so the net increase of the nontrivial-Pauli probability is $\tfrac34-p_0$;
(ii) the denominator $1-e_0$: $p$ is the conditional rate on \emph{non-erased qubits}, and the overall nontrivial-Pauli probability is $(1-e)p+\tfrac34 e$.

Hence $c\le\tfrac34$ if and only if $p_0\ge\tfrac34e_0$. At $e_0=0$ we always have $c=\tfrac34-p_0<\tfrac34$;
but at the work point $(p_0,e_0)=(0.02,0.05)$ suggested in the README, $c=0.768>\tfrac34$.
So ``$R\le3/4$'' is a theorem only in the region $p_0\ge\tfrac34e_0$; in general one should use $c(p_0,e_0)$.
\end{remark}

\begin{corollary}[Integral form: the ``Pauli equivalent'' of erasure]\label{cor:worth}
For every $(p,e)$ define the equivalent $p_{\rm eq}^{(d)}(p,e):=\inf\{p'\in[0,\tfrac34]:\err_d(p',0)\ge\err_d(p,e)\}$. Then
\[
p_{\rm eq}^{(d)}(p,e)\;\le\;p+e\,(\tfrac34-p),
\]
which is Theorem~\ref{thm:partial} with $e_0=0$ and $\delta=e$. If $\err_d(\cdot,0)$ is strictly increasing there is also the lower bound $p_{\rm eq}^{(d)}\ge p$
(monotonicity under superposed erasure, Proposition~\ref{prop:rate-mono}); the Pauli equivalent of the erasure is then sandwiched between $0$ and $e(\tfrac34-p)$,
that is, $0\le R_{\rm eff}\le\tfrac34-p$.\hfill\st{PROVED}
\end{corollary}

\begin{observation}[{Sharpness of the bound: the $[[5,1,3]]$ code attains equality at $e_0=0$}]\label{obs:sharp}
\texttt{theory/checks/exact\_small\_codes.py} enumerates all $4^n$ Pauli errors exactly for three codes with $n\le9$, computes the maximum-likelihood
$\err(p,e)$ in each case, and verifies Theorem~\ref{thm:partial} (all checks pass; see \S\ref{sec:map}). In particular:
\begin{itemize}
\item For the five-qubit code $[[5,1,3]]$ at $e_0=0$ and $p\in\{0.005,0.02,0.06,0.1,0.2,0.3,0.45\}$,
$R^{(d)}=\partial_e\err/\partial_p\err$ and $c=\tfrac34-p$ are \emph{equal} to 8 significant digits.
\item For the Steane code $R^{(d)}/c$ lies between $0.67$ and $0.86$; for the $d=3$ rotated surface code at $p=0.02$, $e_0=0$ one finds $R^{(3)}\approx0.54$ while $c=0.73$.
\end{itemize}
Explanation (unproved): for a perfect code, flagging a single qubit does not change the ML decision on any syndrome, so $\err$ is linear along the segment ``replace the law of qubit $j$ from
$\Dep_p$ by $\Dep_{3/4}$'', and the first-order coefficient is exactly $(\tfrac34-p)\partial_{p_j}\err$.
This shows that the inequality $\partial_e\err_d\le c\,\partial_p\err_d$ cannot be improved uniformly over all codes; the strict gap on the surface code is \emph{code dependent}.
\hfill\st{NUMERICAL}
\end{observation}

\subsection{Complete characterization of the local-free-operation preorder on the code-capacity class}

\begin{definition}[Local free operations]\label{def:local}
On the i.i.d.\ class, a \emph{local free operation} acts qubit by qubit independently: for each qubit the processor draws $(n,f')$ from a Markov kernel
$Q(n,f'\mid f)$ (with $f\in\{0,1\}$ the input flag, $n\in\F^2$ the superposed Pauli, and $f'\in\{0,1\}$ the output flag),
and sets $E'_j=E_j+n$, $F'_j=f'$. ($n$ may depend on $f$, $f'$, and the processor's randomness, but not on $E_j$; this is the general local form of (F1$'$), (F2), and (F2$'$).)
\end{definition}

\begin{theorem}[Preorder theorem: the Pauli-plus-erasure class]\label{thm:preorder}
Let $p,p'\in[0,\tfrac34)$ and $e,e'\in[0,1)$. The following are equivalent:
\begin{enumerate}[(i)]
\item there is a local free operation turning $\Nn_{\rm cc}(p,e)$ into $\Nn_{\rm cc}(p',e')$;
\item either $e'\ge e$ and $\lambda(p')\le\lambda(p)$ (that is, $p'\ge p$); or $e'<e$ and $\mu(p',e')\le\mu(p,e)$,
that is, $(1-e')\lambda(p')\le(1-e)\lambda(p)$.
\end{enumerate}
In particular, (i) implies $\err_d(p',e')\ge\err_d(p,e)$ for every code and every $d$.\hfill\st{PROVED}
\end{theorem}
\begin{proof}
\emph{(ii)$\Rightarrow$(i).} Case $e'\ge e$, $p'\ge p$: first superpose erasure by Lemma~\ref{lem:basic-ops}(b) to take $e\to e'$,
then superpose depolarizing noise by (a) to take $\lambda\to\lambda(p')$ (this needs $\lambda(q)=\lambda(p')/\lambda(p)\in[0,1]$).
Case $e'<e$, $\mu'\le\mu$: by (c) take $\kappa=1-e'/e$, which gives erasure rate $e'$ and
$\lambda_c=(1-e)\lambda(p)/(1-e')$. The condition $\mu'\le\mu$ is exactly $\lambda(p')\le\lambda_c$, and then (a) superposes the remaining depolarizing noise.

\emph{(i)$\Rightarrow$(ii).} For $P\in\F^2\setminus\{0\}$ write $\hat\nu(P)=\sum_a\nu(a)(-1)^{\langle P,a\rangle}$.
Input: with flag 1 (probability $e$) $E$ is uniform and $\hat E(P)=0$; with flag 0 (probability $1-e$) $\hat E(P)=\lambda:=\lambda(p)$.
Let $q_{f,f'}(n)=Q(n,f'\mid f)$ (a sub-probability measure) with total mass $q_{f,f'}$. The Fourier coefficient of the sub-measure with output flag $1$ is
$(1-e)\lambda\,\hat q_{0,1}(P)$ (the part coming from input flag 1 is uniform and has no nonzero modes).
The output class requires ``$E'$ uniform given flag 1'', so $\lambda\hat q_{0,1}(P)=0$ for all $P\ne0$; since $\lambda>0$, $q_{0,1}$ is a uniform measure.
Let $s_0=\sum_n q_{0,0}(n)$; then $\sum q_{0,1}=1-s_0$, and the probability of output flag $0$ is
\[
m':=1-e'=e\,y+(1-e)s_0,\qquad y:=\textstyle\sum_nq_{1,0}(n)\in[0,1].
\]
The Fourier coefficient at $P\neq0$ of the sub-measure with output flag 0 is $(1-e)\lambda\,\hat q_{0,0}(P)$, and the output class requires it to equal
$m'\lambda'$ (with $\lambda'=\lambda(p')$). Since $|\hat q_{0,0}(P)|\le s_0$,
\[
m'\lambda'\le(1-e)\lambda\,s_0.\tag{$*$}
\]
If $e'\ge e$ (that is, $m'\le1-e$): $m'\ge(1-e)s_0$, and substituting into $(*)$ gives $\lambda'\le\lambda$.
If $e'<e$: $s_0\le1$, and $(*)$ gives $(1-e')\lambda'\le(1-e)\lambda$.

The last sentence follows from Theorem~\ref{thm:monotone}.
\end{proof}

\paragraph{Level-1 form.}
For the uncoded qubit, (i) and (ii) are also equivalent to the existence of an arbitrary simulable reduction (label permutations beyond translations) and to $\Phi(P')\le\Phi(P)$ for all $\varphi\in\Phi_4$ (Proposition~\ref{prop:local-level1}); the complete family of monotones is described in Theorem~\ref{thm:complete}.

\begin{remark}[Structure]
(1) Complete monotones: the local free order is completely determined by two functionals: the total unflagged eigenvalue mass $\mu=(1-e)\lambda$ (nonincreasing under any local free operation)
and ``the conditional eigenvalue $\lambda$ is nonincreasing when the erasure rate does not decrease''.
(2) Slope of the boundary of the reachable set: the steepest direction for lowering $e'$ starting from $(p,e)$ is $\mu'=\mu$, whose tangent slope is exactly
$dp'/de'=-c(p',e')$, consistent with the $c$ of Theorem~\ref{thm:rate}: $c$ is the boundary slope of the free order.
(3) The order is not total; for example, if $(p,e)$ and $(p',e')$ satisfy $e'<e$ but $\mu'>\mu$, neither direction is reachable.
\end{remark}

\subsection{Hellinger/Chernoff monotones}

\begin{proposition}[Hellinger monotones]\label{prop:hellinger}
For a noise instance $P=P_{EF}$, $G\in V$, and $s\in(0,1)$ let
$\BC_s(P;G)=\sum_{E,F}P(E,F)^{1-s}P(E+G,F)^s$.
If $P'$ is obtained by a free operation of type (F1) or (F2)/(F2$'$), then $\BC_s(P';G)\ge\BC_s(P;G)$ (for the same $G$).\hfill\st{PROVED}
\end{proposition}
\begin{proof}
These operations are Markov kernels $W$ acting on $(E,F)$ (discarding/coarse-graining flags; superposition of an $E$-independent $N$ is a convolution),
and translation commutes with convolution: $(WP)(\cdot+G)=W(P(\cdot+G))$. The function $(x,y)\mapsto x^{1-s}y^s$ is positively homogeneous and concave, hence superadditive,
$\sum_u(\sum_vW(u|v)x_v)^{1-s}(\sum_vW(u|v)y_v)^s\ge\sum_v x_v^{1-s}y_v^s$. Take $x=P$, $y=P(\cdot+G)$.
\end{proof}

For $\Nn_{\rm cc}(p,e)$ and any nonzero single-qubit shift $G_j\in\{X,Y,Z\}$ (the three are symmetric), the single-qubit coefficient is
\[
b_s(p,e)=e+(1-e)\Big[(1-p)^{1-s}(\tfrac p3)^s+(\tfrac p3)^{1-s}(1-p)^s+\tfrac{2p}3\Big],
\]
where the first term comes from $F=1$ (the uniform distribution is invariant under the shift). Since $b_s=b_{1-s}$ and $b_s$ is convex in $s$, $\min_s b_s=b_{1/2}$. Put
\[
\beta_{\rm dep}(p)=2\sqrt{p(1-p)/3}+2p/3,\qquad B(p,e)=b_{1/2}(p,e)=e+(1-e)\beta_{\rm dep}(p).
\]
$B$ is the \emph{single-qubit Bhattacharyya parameter}. For each $s$, the same argument as in Theorem~\ref{thm:partial} together with Proposition~\ref{prop:hellinger}
(partial flag discarding does not decrease $b_s$) gives $\partial_eb_s\le c\,\partial_pb_s$, that is, the Hellinger exchange rate satisfies
\[
R_s(p,e)=\frac{\partial_eb_s}{\partial_pb_s}\le c(p,e).
\]
Write $R_B:=R_{1/2}=\dfrac{1-\beta_{\rm dep}(p)}{(1-e)\,\beta_{\rm dep}'(p)}$.

\subsection{The gap problem}\label{sec:gap}

The free order $\succeq_{\rm free}$ is a sub-order of the \emph{operational order} (comparing $\err_d$ or $\alpha$ for a given code). Let $R_\alpha$ be the true exchange rate on the surface code.
\begin{problem}[Exchange-rate gap]\label{prob:gap}
Quantify $\Delta:=c-R_\alpha\ge0$, and decide whether the gap comes from
\begin{enumerate}[(A)]
\item \emph{code dependence}: the \emph{code-universal} operational order (requiring $\err$ not to increase for all linear decoding structures and all $n$) coincides with the local free order,
and the $\alpha$-order of the surface code is merely a coarser order; or
\item \emph{incompleteness of the free operations}: there are free operations outside Definition~\ref{def:local} (nonlocal or code dependent) that would make the order finer?
\end{enumerate}
\end{problem}

\begin{conjecture}[Code-universal order = local free order]\label{conj:universal}
Within $\Nn_{\rm cc}(p,e)$, the \emph{code-universal finite-$n$ operational order} coincides with the local free order of Theorem~\ref{thm:preorder}.
\hfill\st{CONJECTURE}
\end{conjecture}
\noindent\emph{Reformulation (constitutional audit).} The all-losses version of this conjecture (compare $\Phi(P')\le\Phi(P)$ for all $\varphi\in\Phi_\Lambda$, over all structures) is proved by taking the uncoded qubit (Proposition~\ref{prop:local-level1}, Remark~\ref{rem:conj-reform}); what remains open is the single-loss statement for $\err=1-\Phi_{\max}$.

\noindent\emph{Evidence and route.} Observation~\ref{obs:sharp} gives first-order evidence at $e_0=0$: the boundary slope $c$ is attained by some code.
Route to a proof: (S1) for each $n$ and linear structure $(\sigma,L)$, the comparison of $\err$ is a Blackwell comparison of the ``experiments'' given by $n$ independent copies
(Blackwell--Sherman--Stein: experiment $A$ is at least as good as $B$ for all priors and losses if and only if $B$ is a randomized processing of $A$ \cite{Blackwell});
(S2) the decision problems arising from linear structures form a proper subclass of all decision problems, and one must show that they suffice to separate the two functionals $\mu$ and $(e,\lambda)$ of the local free order;
(S3) the asymptotic version is characterized, via the large-sample Blackwell order, by the family of all R\'enyi divergences \cite{MPST}, which corresponds exactly to $\{b_s\}_{s\in(0,1)}$ and its higher-order generalizations.
If the conjecture holds, the gap is case (A) and $c-R_\alpha$ measures information specific to the surface code rather than a deficiency of the free operations;
if it is refuted by a counterexample from some code (for example one other than $[[5,1,3]]$), new free operations are needed (case (B)).

\begin{conjecture}[Completeness of $B$: hypothesis $H_B$]\label{conj:HB}
In some neighborhood within $\Nn_{\rm cc}$, the $\alpha(p,e)$ of the surface code is a function of $B(p,e)$, so that $R_\alpha=R_B$.
\hfill\st{CONJECTURE}
\end{conjecture}
\noindent This is testable: the numerics of P1 give $R_\alpha$, to be compared with the analytic $R_B$. $R_\alpha\neq R_B$ would show that
the Bhattacharyya monotone is incomplete, and the next candidates are $b_s$ ($s\ne\tfrac12$). Lemma~\ref{lem:union} (upper bound) and Proposition~\ref{prop:genie}
(lower bound) already sandwich $\alpha$ between two functions depending only on $B$, but a sandwich does not constrain derivatives, so $H_B$ needs its own test.

\subsection{Analytic formula for derivatives and the envelope argument}\label{sec:envelope}

The sampling scheme of P1 (README \S2) obtains $\partial_\theta\pL$ directly from one table $f(k,w)$, without decoding at neighboring parameter points. For decoders independent of $(p,e)$ (MWPM) this is an identity;
for the maximum-likelihood decoder it relies on the following argument.

\begin{remark}[Envelope argument for ML derivatives]\label{rem:envelope}
Let $D^\ast$ be the ML decoder at $\theta^\ast=(p^\ast,e^\ast)$ (a fixed map). For a fixed decoder $D$,
$\pL(\theta;D)=\sum_{k,w}P_\theta(k,w)\,f_D(k,w)$, where $f_D$ does not depend on $\theta$. Optimality gives $\err(\theta)=\min_D\pL(\theta;D)\le\pL(\theta;D^\ast)$,
with equality at $\theta^\ast$. The difference $g(\theta)=\pL(\theta;D^\ast)-\err(\theta)\ge0$ satisfies $g(\theta^\ast)=0$, so wherever both are differentiable at $\theta^\ast$ we have $\nabla g(\theta^\ast)=0$, that is,
\[
\partial_\theta\err\big|_{\theta^\ast}=\sum_{k,w}\partial_\theta P_\theta(k,w)\,f_{D^\ast}(k,w).
\]
The ML decoder depends on $\theta$ only through $p$: the erased positions are known, so $e$ never enters the posterior given the flags. Therefore
the $e$-derivative is an exact identity for ML, and the envelope argument is needed only for $\partial_p$.
Hence $\partial_\theta\ln\pL=\langle s_\theta\rangle_{\rm fail}$ (with $s_\theta=\partial_\theta\ln P_\theta$ the score function and the conditional expectation taken over failure events),
and $\partial\alpha/\partial\theta$ is minus the least-squares slope of $\langle s_\theta\rangle_{\rm fail}$ against $d$ (the least-squares slope is linear in the data and commutes with differentiation).
Differentiability: at finite $d$, $\err$ is the minimum of finitely many smooth functions and is differentiable almost everywhere; kinks may occur at ties between decisions.
\hfill\st{SKETCH}
\end{remark}

\begin{observation}[Exact verification of the envelope identity]\label{obs:envelope}
Check C8 of \texttt{exact\_small\_codes.py}: for $[[5,1,3]]$, the Steane code, and the $d=3$ rotated surface code, at $(p^\ast,e^\ast)\in\{(0.04,0.05),(0.1,0.2)\}$,
the derivatives $\partial_p\err,\partial_e\err$ of the exact ML agree with the derivatives of the failure rate of the \emph{frozen decoder} to within relative $2\times10^{-10}$.
In addition, \texttt{experiments/p1\_erasure\_pauli/envelope\_check.py} tests the $p$-derivative at $d=5,7$ with the exact transfer-matrix ML decoder on common random samples
of each stratum $(k,w)$. A central difference has a deterministic $O(h^2)$ discretisation error (about $+0.3\%$ at $d=5$ and $+1.9\%$ at $d=7$ for $h=0.01$, $p=0.06$, growing four times for $2h$),
so the test compares the finite difference with re-matched decoders ($B_{\rm matched}$) against the \emph{same} finite difference with the frozen decoder ($B_{\rm frozen}$).
Their difference is at most $0.04\%$ of the derivative (within $1.6$ bootstrap standard errors) in four cases, one of which has erasure strata; one deviation of $2.05$ standard errors ($0.1\%$ of the derivative) occurs at step $2h$, $d=7$.
The identity has not yet been tested directly for the truncated-MPS ML decoder or at $d\ge9$; the MPS decoder's decisions agree with the exact decoder on all tested samples (README, M1), so it is expected to transfer.\hfill\st{NUMERICAL}
\end{observation}

\section{条件认证:有限样本上界}\label{sec:cert}

目标:在模型类内,\emph{直接}对大码距的 $\pL(d)$ 给出有限样本上界,而不是先估计 $\alpha$ 再外推。
证明骨架分三层:\emph{(A) 统计层}(由数据得到参数的置信上界);\emph{(B) 单调层}(参数的上界给出模型的``支配'',
由定理~\ref{thm:monotone} 得 $\err$ 的上界);\emph{(C) 计算层}(对支配模型给出 $\err_d$ 的显式上界)。
本节完整证明码容量情形(定理~\ref{thm:cc-cert}),并把含突发与芯片事件的情形(定理~\ref{thm:burst-cert})
化为三个明确的开放引理。

\paragraph{关于“哪个解码器”。}
本节所有上界都是对\emph{最大似然(Bayes 最优)解码}的上界,表述为“存在解码器(已知噪声参数下的 ML)
使逻辑失败概率不超过 $U$”。由最优性,它\emph{不是}对实验中实际使用的次优解码器(如 MWPM)失败率的上界。
反之,本节的\emph{下界}(地板,引理~\ref{lem:floor})对任何解码器成立。

\subsection{(C) 联合(Bhattacharyya)上界}

\begin{lemma}[联合上界]\label{lem:union}
设 $(E,F)$ 为线性解码结构 $(V,\sigma,L)$ 上的任意噪声实例,令 $K=\ker\sigma\setminus\ker L$
(不被检测却翻转可观测量的故障)。则
\[
\err(E,F)\;\le\;\tfrac12\sum_{G\in K}\BC(P;G),\qquad
\BC(P;G):=\sum_{E,F}\sqrt{P(E,F)\,P(E+G,F)}.
\]
\hfill\st{PROVED}
\end{lemma}
\begin{proof}
固定观测 $y=(s,f)$,记 $P(\ell,y)=\Prob[L=\ell,Y=y]$,ML 判决为 $\ell^\ast(y)$。则
$\err=\sum_y\sum_{\ell\ne\ell^\ast}P(\ell,y)$。对 $\ell\ne\ell^\ast$ 有 $P(\ell,y)\le P(\ell^\ast,y)$,故
$P(\ell,y)\le\sqrt{P(\ell,y)P(\ell^\ast,y)}$,于是
\[
\err\le\sum_y\sum_{\{\ell,\ell'\}}\sqrt{P(\ell,y)P(\ell',y)}\quad(\text{对无序标签对求和}).
\]
记 $c(s,\ell)=\{E:\sigma(E)=s,L(E)=\ell\}$,有 $P(\ell,s,f)=\sum_{A\in c(s,\ell)}P(A,f)$。
由 $\sqrt{\sum_ix_i}\le\sum_i\sqrt{x_i}$ 得
\[
\sqrt{P(\ell,s,f)P(\ell',s,f)}\le\sum_{A\in c(s,\ell)}\sum_{B\in c(s,\ell')}\sqrt{P(A,f)P(B,f)}.
\]
对 $s,\{\ell,\ell'\},f$ 求和,右边变成对满足 $\sigma(A)=\sigma(B)$、$L(A)\ne L(B)$ 的\emph{无序}故障对 $\{A,B\}$ 求和。
令 $G=A+B\in K$,$(A,B)\mapsto(A,G)$ 是有序对到 $V\times K$ 的双射,无序对数为其一半,
故得 $\tfrac12\sum_{G\in K}\sum_{A,f}\sqrt{P(A,f)P(A+G,f)}$。
\end{proof}

\begin{corollary}[码容量的积形式与重量枚举]\label{cor:product}
对 $\Nn_{\rm cc}(p,e)$ 与稳定子码 $(n,S)$,$K=N(S)\setminus S$,且对 $G\in K$
$\BC(P;G)=B(p,e)^{\wt(G)}$,其中 $B=e+(1-e)\beta_{\rm dep}(p)$($\wt$ 为泡利重量)。故
\[
\err_d(p,e)\le\tfrac12\widetilde W_d\big(B(p,e)\big),\qquad
\widetilde W_d(x)=\sum_{G\in N(S)\setminus S}x^{\wt(G)}=W_{N(S)}(x)-W_S(x),
\]
并由 MacWilliams 恒等式 $W_{N(S)}(x)=\dfrac1{|S|}\sum_{g\in S}(1+3x)^{n-\wt(g)}(1-x)^{\wt(g)}$ 可由稳定子群的重量分布计算。
\hfill\st{PROVED}
\end{corollary}
\begin{proof}
单比特:标记 $f=1$(概率 $e$)时分布均匀,平移不变,贡献 $e\cdot1$;标记 $f=0$ 时贡献
$(1-e)\sum_a\sqrt{\Dep_p(a)\Dep_p(a+G_j)}=(1-e)\beta_{\rm dep}(p)$($G_j\in\{X,Y,Z\}$ 三者对称;
配对 $(I,G_j),(G_j,I)$ 给出 $2\sqrt{(1-p)p/3}$,另外两个非平凡泡利互相配对给出 $2p/3$)。
$G_j=I$ 的比特贡献 $1$。独立性给出积形式。MacWilliams 恒等式为标准结果。
\end{proof}

\noindent 对 $d=3$ 的三个码,不等式 $\err\le\tfrac12\widetilde W(B)$ 在 $\texttt{theory/checks}$ 中全部通过(检验 C5)。
一般 $d$ 时 $W_S$ 可由张量网络计算(每格点 4 维腿,转移矩阵宽度随 $d$ 指数增长),具体可行的 $d$ 范围待 M1 之后基准测试,\emph{此处不作承诺}。

\subsection{(A) 统计层:参数的置信上界}

\begin{lemma}[由标记估计 $e$]\label{lem:flags}
$M$ 个独立样本中共 $nM$ 个比特-样本,标记 $F_{ij}$ i.i.d.\ Bernoulli($e$)。令
$\bar e=\bar e(\delta_1)$ 为 Clopper--Pearson 单侧上限,则 $\Prob[e\le\bar e]\ge1-\delta_1$。\hfill\st{PROVED}
\end{lemma}

\begin{lemma}[由综合征频率估计 $\lambda$]\label{lem:lambda}
取一个固定的 $X$ 型稳定子生成元 $g_0$,重量 $w_0$。对其支撑上没有任何标记的那些样本(记其个数为 $M'$),
综合征比特为 i.i.d.\ Bernoulli,参数 $\tfrac12(1-\lambda^{w_0})$(以标记模式为条件,精确成立)。
令 $\overline{P}$ 为该参数的 Clopper--Pearson 上限(置信水平 $1-\delta_2$),则
$\lambda\ge\underline\lambda:=(1-2\overline P)_+^{1/w_0}$,即 $p\le\bar p:=\tfrac34(1-\underline\lambda)$,
概率 $\ge1-\delta_2$。\hfill\st{PROVED}
\end{lemma}
\begin{proof}
没有标记时 $g_0$ 支撑上每个比特以概率 $p$ 去极化;$X$ 型检测器翻转当且仅当支撑上 $\{Y,Z\}$ 型错误的个数为奇数,
单比特贡献本征值 $1-\tfrac43p=\lambda$,故 $\Pr[\text{bit}=1]=\tfrac12(1-\lambda^{w_0})$。
以标记为条件时各样本独立,故为精确二项。
\end{proof}

\begin{theorem}[码容量情形的有限样本认证]\label{thm:cc-cert}
设数据为码容量噪声 $\Nn_{\rm cc}(p,e)$ 的 $M$ 个独立样本(标记与综合征),$\delta=\delta_1+\delta_2$。
令 $\bar e,\bar p$ 如引理~\ref{lem:flags}、\ref{lem:lambda}。则以概率至少 $1-\delta$,对\emph{一切} $d$,
\[
\err_d(p,e)\;\le\;\err_d(\bar p,\bar e)\;\le\;\tfrac12\widetilde W_d\big(B(\bar p,\bar e)\big).
\]
\hfill\st{PROVED}
\end{theorem}
\begin{proof}
在事件 $\{e\le\bar e\}\cap\{p\le\bar p\}$ 上,由定理~\ref{thm:preorder}(情形 $e'\ge e$、$p'\ge p$)
$\Nn_{\rm cc}(\bar p,\bar e)$ 可由 $\Nn_{\rm cc}(p,e)$ 经局域自由操作得到(叠加擦除,再叠加去极化),
故定理~\ref{thm:monotone} 给出第一个不等式;第二个是推论~\ref{cor:product}。
该事件的概率由并集界 $\ge1-\delta_1-\delta_2$。
\end{proof}

这是一个\emph{不依赖任何 $\pL$ 数据}的认证:标记和综合征统计足以给出所有 $d$ 的上界。代价是联合上界的松弛
(它的阈值远低于真实 ML 阈值)。把联合上界换成更紧的、同样单调的上界(例如用测得的小 $d$ 失败率校准)是下一步。

\subsection{含局域突发与芯片尺度事件的情形}

设模型类为定义~\ref{def:events}。记 $q_c=1-e^{-r_c\,T_{\rm tot}}$ 为总时长 $T_{\rm tot}$ 内至少发生一次芯片尺度事件的概率。

\begin{lemma}[芯片尺度事件的地板]\label{lem:floor}
令 ${\err}^{(0)}$ 为去掉芯片尺度层后的同一噪声的 $\err$。则
\[
\err\;\le\;q_c+(1-q_c)\,{\err}^{(0)}.
\]
若进一步假设 \textbf{(K)}:给定发生芯片事件,可观测量 $L(E)\in\F^q$ 均匀且与 $Y=(\sigma(E),F)$ 独立,则
\[
\err\;\ge\;q_c\,(1-2^{-q}).
\]
于是 $\err\sim q_c\propto r_cT_{\rm tot}$ 的“地板”与 $d$ 无关;上界对 (K) 不需要,下界对任何解码器的失败率同样成立。
\hfill\st{PROVED}
\end{lemma}
\begin{proof}
上界:用去掉芯片层后最优的 ML 解码器 $\hat g_0$ 处理完整问题。设 $C$ 为“至少一个芯片事件”。
$\Prob[\hat g_0\text{ 失败}]\le\Prob[C]+\Prob[C^c]\,{\err}^{(0)}$,因为在 $C^c$ 上数据服从去掉芯片层的噪声(各层独立)。
$\err$ 不超过任何解码器的失败率。

下界:考虑知道 $C$ 的先知解码器,其失败率不超过 $\err$(定理~\ref{thm:monotone} 的 (F1) 情形:真实观测是先知观测的粗粒化)。
在 $C$ 上由 (K) 任何猜测的失败概率为 $1-2^{-q}$。
(K) 对“全芯片完全去极化事件”成立:均匀的 $E_{\rm chip}$ 使 $(\sigma,L)(E)$ 在 $\F^{n-k}\times\F^{2k}$ 上均匀(满射),
故 $L$ 与 $\sigma$ 独立;标记独立于 $E_{\rm chip}$。
\end{proof}

\begin{lemma}[突发混合的联合上界]\label{lem:burst-union}
设 $E=E_{\rm bg}+E_{\rm burst}$,以突发构型 $\mathcal C$ 为条件时噪声是比特间独立的积测度,构型分布为 $\pi$。则
\[
\BC(P;G)\;\le\;\sum_{\mathcal C,\mathcal C'}\sqrt{\pi(\mathcal C)\pi(\mathcal C')}\;\prod_{j}\beta_j(\mathcal C,\mathcal C';G_j),
\]
其中 $\beta_j(\mathcal C,\mathcal C';G_j)=\sum_{a,f}\sqrt{P_{\mathcal C,j}(a,f)P_{\mathcal C',j}(a+G_j,f)}$。
\hfill\st{PROVED}
\end{lemma}
\begin{proof}
$P=\sum_{\mathcal C}\pi(\mathcal C)P_{\mathcal C}$。由 $\sqrt{\sum x_i}\le\sum\sqrt{x_i}$,
\[
\sqrt{P(E,F)P(E+G,F)}\le\sum_{\mathcal C,\mathcal C'}\sqrt{\pi(\mathcal C)\pi(\mathcal C')}\sqrt{P_{\mathcal C}(E,F)P_{\mathcal C'}(E+G,F)}.
\]
对 $(E,F)$ 求和并用积测度。
\end{proof}
\noindent 注意不能简单地“以 $\mathcal C$ 为条件、把 $\mathcal C$ 当作已知”:那样得到的是知道构型的解码器的失败率,
它\emph{低于} $\err$(先知更强),不是上界。引理~\ref{lem:burst-union} 是对未知构型的正确处理,但右端把两条独立的突发构型同时求和,
结果偏松。

\begin{theorem}[含突发与芯片事件的条件认证(目标陈述)]\label{thm:burst-cert}
设真实参数 $\theta=(p,e,r,\vartheta,r_c)$,形状 $\vartheta=(R,T_b,p_b)$ 属于已知的有限类 $\Theta_\vartheta$(若问题~\ref{prob:O1} 的单调性成立,可代以形状上限;否则下式右端对 $\vartheta\in\Theta_\vartheta$ 取最大值)。
假设由数据得到的置信上限满足 $\Prob[p\le\bar p,\ e\le\bar e,\ r\le\bar r,\ r_c\le\bar r_c]\ge1-\delta$。则
以概率 $\ge1-\delta$,对一切 $d,T_{\rm tot}$,
\[
\err_{d,T_{\rm tot}}(\theta)\;\le\;\bar q_c+(1-\bar q_c)\,\tfrac12\sum_{G\in K}\ \sum_{\mathcal C,\mathcal C'}\sqrt{\pi_{\bar\theta}(\mathcal C)\pi_{\bar\theta}(\mathcal C')}\prod_j\beta_j(\mathcal C,\mathcal C';G_j),
\]
其中 $\bar q_c=1-e^{-\bar r_cT_{\rm tot}}$,$\pi_{\bar\theta}$ 为参数上限对应的突发构型分布。\hfill\st{PLAN}
\end{theorem}

\noindent\textbf{已证明的组成部分:}速率方向的单调性(命题~\ref{prop:rate-mono}),地板(引理~\ref{lem:floor}),
联合上界(引理~\ref{lem:union}、\ref{lem:burst-union});把 $\bar\theta$ 代入的论证与定理~\ref{thm:cc-cert} 相同。
\textbf{尚缺三块:}
\begin{problem}[O1:形状参数的单调性]\label{prob:O1}
无标记突发下,$\err$ 对 $p_b,R,T_b$ 是否单调?(注~\ref{rem:shape} 说明自由操作不给出;若不成立,
则须对形状取有限类上的最坏情形,或对形状作联合上界。)
\end{problem}
\begin{problem}[O2:污染下的参数识别]\label{prob:O2}
突发与芯片事件污染综合征统计,引理~\ref{lem:lambda} 的 i.i.d.\ 前提失效。需要对 $p$ 的稳健置信上限(例如排除高权重检测事件窗口后的估计)
与对 $r,r_c$ 的置信上限(由事件检测统计,与 P2 的泊松下限 §\ref{sec:design} 衔接)。
\end{problem}
\begin{problem}[O3:突发和的显式估计]\label{prob:O3}
把引理~\ref{lem:burst-union} 右端的二重构型和化为显式的指数函数(团簇展开);需给出收敛域,
并说明在 $d\gg R$ 时它只改变指数、不产生与 $d$ 无关的地板。
\end{problem}

\paragraph{数值检验的分工(见 §\ref{sec:map})。}
电路级 stim + MWPM 的 P2 数据能检验:(i) 地板 $\err\ge q_c(1-2^{-q})$ 与 $\propto r_cT_{\rm tot}$ 的标度
(下界对任何解码器成立,MWPM 的失败率不得低于它);(ii) 参数估计器的无偏性(M3 验收项);
(iii) 突发使指数变小而不是产生地板。\emph{上界}部分针对 ML,MWPM 数据不能证伪它;
须用小 $d$ 的精确 ML 或码容量下带突发簇的张量网络 ML(M1)检验。

\section{Information-theoretic bounds and experimental design}\label{sec:info}

\subsection{A sandwich for $\alpha$: the Chernoff upper bound and the union lower bound}

\begin{proposition}[Chernoff upper bound, with an explicit finite-$d$ form]\label{prop:genie}
Let $d\ge3$ be odd, consider the rotated surface-code family (conventions of \texttt{lcd.codes.RotatedSurfaceCode}) and noise $\Nn_{\rm cc}(p,e)$ with $0\le p\le\tfrac34$, $0\le e\le1$, $(p,e)\ne(0,0)$. Then
\[
\alpha_+(p,e)\;\le\;\ln\frac1{B(p,e)},\qquad B(p,e)=e+(1-e)\Big(2\sqrt{\tfrac{p(1-p)}3}+\tfrac{2p}3\Big).
\]
More precisely, for every $d$ we have $\err_d\ge\varepsilon_{\rm genie}(d)$, where
$\varepsilon_{\rm genie}(d)=\tfrac12\bigl(1-\mathrm{TV}(P_1^{\otimes d},Q_1^{\otimes d})\bigr)$,
$P_1$ is the joint law of a single qubit's $(a,f)$ and $Q_1=P_1(\cdot+X)$ is the shifted law, and
\[
\varepsilon_{\rm genie}(d)\;\ge\;\tfrac12\,B^d\exp\!\Big(-\tfrac12\sqrt{d\,\sigma^2}\Big),\qquad
\sigma^2=\frac{2(1-e)\sqrt{(1-p)p/3}}{B}\,\ln^2\!\frac{3(1-p)}{p}\quad(\sigma^2=0\text{ if }p=0).
\]
\hfill\st{PROVED}
\end{proposition}
\begin{proof}
Take the $d$ qubits $T$ of an interior column $c_0$, $1\le c_0\le d-2$, and let $\bar X_T$ be the $X$-type operator on this column.

\emph{Step 1: $\ker\sigma\cap\mathcal P_T=\{I,\bar X_T\}$, and $\bar X_T\notin S$.} A plaquette with top-left corner $(r,c)$ covers $(r,c),(r,c+1),(r+1,c),(r+1,c+1)$ and is of $Z$ type iff $r+c$ is odd.
For each $r\in\{0,\dots,d-2\}$ the two bulk plaquettes with corners $(r,c_0-1)$ and $(r,c_0)$ exist (because $1\le c_0\le d-2$) and have opposite types; each meets column $c_0$ exactly in $\{(r,c_0),(r+1,c_0)\}$.
Let $(x,z)$ be supported on $T$ with zero syndrome. The $Z$-type plaquette of the pair forces $x_r=x_{r+1}$ for every $r$, so $x\in\{0,\mathbf 1_T\}$; the $X$-type one forces $z_r=z_{r+1}$, so $z$ is constant.
The weight-$2$ $X$-type checks on the top edge have corners $(-1,c)$ with $c$ odd and cover $(0,c),(0,c+1)$; exactly one of them ($c=c_0$ if $c_0$ is odd, $c=c_0-1$ if $c_0$ is even) contains $(0,c_0)$, and it meets $T$ only there. Hence $z_0=0$ and $z=0$.
Conversely $\bar X_T$ has zero syndrome: bulk $Z$ plaquettes meet the column in $0$ or $2$ qubits and the $Z$-type boundary checks lie on columns $0$ and $d-1$. It anticommutes with the $Z$-type logical on row $0$ (overlap $1$), so $\bar X_T\notin S$ and it flips the label.
(Check K1 of \texttt{checks/completion\_checks.py} verifies the statement over $\F$ for every interior column, $d=3,\dots,41$.)

\emph{Step 2: the genie.} The genie knows $F$, $E_{T^c}$, and $\sigma(E)$, which is equivalent to knowing $\sigma(E_T)$. By Step 1, $E_T$ can only be $a$ or $a+\bar X_T$,
and the two have different labels. The genie's Bayes error is $\tfrac12\sum_{a,f_T}\min\big(P_T(a,f_T),P_T(a+\bar X_T,f_T)\big)=\varepsilon_{\rm genie}$.
The real observation $(\sigma(E),F)$ is a coarse-graining (F1) of the genie's observation, so $\err_d\ge\varepsilon_{\rm genie}$ (Theorem~\ref{thm:monotone}).

\emph{Step 3: the explicit bound.} Write $P=P_1^{\otimes d}$, $Q=Q_1^{\otimes d}$ and $\Lambda=\ln(P/Q)$ where $PQ>0$. Since $\min(P,Q)=\sqrt{PQ}\,e^{-|\Lambda|/2}$ (both sides vanish where $PQ=0$) and $\sum_x\sqrt{P(x)Q(x)}=B^d$
($B$ is the Bhattacharyya coefficient of $P_1$ and $Q_1$, \S\ref{sec:exchange}),
\[
\varepsilon_{\rm genie}=\tfrac12B^d\,\E_R\big[e^{-|\Lambda|/2}\big],\qquad R:=\sqrt{PQ}/B^d=R_1^{\otimes d},\quad R_1=\sqrt{P_1Q_1}/B .
\]
Under $R$, $\Lambda=\sum_i\lambda(x_i)$ with $x_i$ i.i.d.\ $\sim R_1$. The involution $x\mapsto x+X$ exchanges $P_1$ and $Q_1$, so it preserves $R_1$ and maps $\lambda$ to $-\lambda$; hence $\E_{R_1}\lambda=0$ and $\E_R\Lambda^2=d\,\E_{R_1}\lambda^2=d\sigma^2$.
On an unflagged qubit $\lambda=\pm\ln\frac{3(1-p)}p$ for $a\in\{I,X\}$ (each of $R_1$-mass $(1-e)\sqrt{(1-p)p/3}/B$) and $\lambda=0$ for $a\in\{Y,Z\}$; on a flagged qubit $\lambda=0$; this gives the stated $\sigma^2$.
By Jensen's inequality (convexity of $\exp$) and Cauchy--Schwarz, $\E_R e^{-|\Lambda|/2}\ge e^{-\E_R|\Lambda|/2}\ge e^{-\sqrt{\E_R\Lambda^2}/2}=e^{-\sqrt{d\sigma^2}/2}$.
Finally $-\tfrac1d\ln\err_d\le\ln\tfrac1B+\tfrac{\ln2}d+\tfrac{\sigma}{2\sqrt d}$, whose $\limsup$ is $\ln(1/B)$.
Note that the flags live in the joint alphabet $(a,f)$ and the genie sees them, so at positions with $F=1$ the two hypotheses are \emph{identical},
contributing $e\cdot1$ to $B$; this is the origin of the $e$ term in $B$.
\end{proof}

\noindent The finite-$d$ bound replaces the earlier appeal to the asymptotic Chernoff theorem \cite[\S11.9]{CT}; its exponent is the Chernoff information $-\ln B$ (the minimum over $s$ of $b_s$ is at $s=\tfrac12$, \S\ref{sec:exchange}).
Check K2 compares the bound with the exact $\varepsilon_{\rm genie}(d)$ for $d\le61$ (ratio $\ge1.67$ on the tested grid) and the slope of $-\ln\varepsilon_{\rm genie}$ with $\ln(1/B)$ (within $2\%$ between $d=101$ and $201$).
For $d=3$, $\err\ge\varepsilon_{\rm genie}$ has been verified by exact computation (check C6); \texttt{checks/genie\_large\_d.py} tests it against the ML data of P1 for $d=5$--$13$.

\begin{proposition}[Union lower bound]\label{prop:lower}
Let $\Phi(x)=\limsup_{d\to\infty}\tfrac1d\ln\widetilde W_d(x)$ (Corollary~\ref{cor:product}). Then
$\alpha_-(p,e)\ge-\Phi\big(B(p,e)\big)$.\hfill\st{PROVED}
\end{proposition}
\begin{proof}
Corollary~\ref{cor:product}: $\err_d\le\tfrac12\widetilde W_d(B)$; take $-\tfrac1d\ln$ and then $\liminf$.
\end{proof}

\begin{corollary}[Sandwiched between two functions of $B$ only]\label{cor:sandwich}
$-\Phi(B(p,e))\le\alpha_-\le\alpha_+\le-\ln B(p,e)$.\hfill\st{PROVED}
\end{corollary}

The sandwich does not constrain the derivatives of $\alpha$, so it does not directly give upper and lower bounds on the exchange rate; its role is to give the testable hypothesis $H_B$
(Conjecture~\ref{conj:HB}): if $\alpha$ is a function of $B$ in a neighborhood, then $R_\alpha=R_B$.

\subsection{Reference values for the exchange rate}

The table below lists three quantities at the README work points, all computable from closed-form expressions. Here $c$ is the rigorous upper bound (Theorem~\ref{thm:rate});
$R_B$ and $R_{\rm hash}$ are \emph{heuristic references}, \textbf{not} theorems about the $R_\alpha$ of the surface code, and serve only to give the P1 numerics an order-of-magnitude expectation:
\[
R_{\rm hash}(p,e)=\frac{2-H_4(p)}{(1-e)\,\log_2\!\frac{3(1-p)}{p}},\quad
H_4(p)=-(1-p)\log_2(1-p)-p\log_2\tfrac p3,
\]
the slope of the level curves of the quantum hashing capacity $Q_{\rm hash}=1-2e-(1-e)H_4(p)$ (whose zero level curve gives $p\approx18.9\%$ at $e=0$
and $e=50\%$ at $p=0$).

\begin{center}
\begin{tabular}{cc ccc}
\toprule
$p_0$ & $e_0$ & $c=\frac{3/4-p_0}{1-e_0}$ (upper bound) & $R_B$ & $R_{\rm hash}$ \\
\midrule
0.02 & 0    & 0.7300 & 0.1784 & 0.2537 \\
0.02 & 0.02 & 0.7449 & 0.1820 & 0.2589 \\
0.02 & 0.05 & 0.7684 & 0.1877 & 0.2671 \\
0.04 & 0    & 0.7100 & 0.2212 & 0.2746 \\
0.04 & 0.02 & 0.7245 & 0.2257 & 0.2802 \\
0.04 & 0.05 & 0.7474 & 0.2328 & 0.2891 \\
0.06 & 0    & 0.6900 & 0.2444 & 0.2840 \\
0.06 & 0.02 & 0.7041 & 0.2494 & 0.2898 \\
0.06 & 0.05 & 0.7263 & 0.2573 & 0.2989 \\
\bottomrule
\end{tabular}
\end{center}
(These values are printed by check C7 of \texttt{theory/checks/exact\_small\_codes.py}. The same script also gives the finite-size exchange rates $R^{(3)}$ of the $d=3$ codes,
for example $0.54$ for the surface code ($p_0=0.02,e_0=0$); it satisfies $R^{(d)}\le c$ but lies above $R_B$ and $R_{\rm hash}$;
how the finite-size $R^{(d)}$ evolves with $d$ is exactly what P1 is to measure.)

\begin{remark}[Threshold secant]
Using the known thresholds $p_c\approx0.189$ (at $e=0$) and $e_c=0.5$ (at $p=0$), the secant slope of the zero level curve is $0.189/0.5\approx0.38$,
far below $c\approx0.75$. This is only a heuristic comparison of a secant (at threshold $\alpha=0$, and the local derivative at a work point is different), not a theorem.
As a rigorous consistency check: Corollary~\ref{cor:worth} gives $p_{\rm eq}\le3e/4$, so the pure-erasure threshold must satisfy
$e_c\ge\tfrac43p_c\approx0.25$, which is compatible with $e_c=0.5$.
\end{remark}

\subsection{Experimental design and detection limits}\label{sec:design}

This part corresponds to P2. Let $T$ be the total observation time (number of rounds) and let events be a Poisson process of rate $r$.

\begin{proposition}[Poisson lower bound]\label{prop:poisson}
Let $H_0$: no burst events; $H_1$: burst events form a Poisson process of intensity $r$ (of arbitrary shape), with all other layers independent and identical.
For any test $\varphi$ based on all data over the total time $T$, if the type-I error under $H_0$ is $\le a$, then the power under $H_1$ satisfies
\[
\Prob_1[\varphi=1]\;\le\;1-e^{-rT}(1-a).
\]
Therefore reaching power $1-\delta$ requires $T\ge\ln\!\big((1-a)/\delta\big)/r$. No method that exploits space-time patterns can beat the $1/r$ scaling.\hfill\st{PROVED}
\end{proposition}
\begin{proof}
Under $H_1$, with probability $e^{-rT}$ no event occurs at all, and then the data have the same law as under $H_0$. Hence
$\Prob_1[\varphi=1]\le(1-e^{-rT})+e^{-rT}\Prob_0[\varphi=1]\le1-e^{-rT}(1-a)$.
\end{proof}

\begin{proposition}[Chain of exponents for class discrimination]\label{prop:exponents}
Let events carry a label $a\in\mathcal A$ (the class), and let the class intensities under the two hypotheses be $r_a^{(0)},r_a^{(1)}$. With events fully observed,
the Stein exponent and the Chernoff exponent per unit time are
\[
D=\sum_a\Big[r^{(1)}_a\ln\frac{r^{(1)}_a}{r^{(0)}_a}-r^{(1)}_a+r^{(0)}_a\Big],\qquad
C=\max_{s\in[0,1]}\sum_a\Big[s\,r^{(1)}_a+(1-s)\,r^{(0)}_a-(r^{(1)}_a)^s(r^{(0)}_a)^{1-s}\Big]
\]
(the relative entropy and Chernoff information of Poisson processes). If the actual data $Z$ are the output of the event process through a detector channel $W$, then
\[
D_{\rm count}\ \le\ D_{\rm pattern}\ \le\ D_{\rm event}=D,\qquad C_{\rm count}\le C_{\rm pattern}\le C,
\]
where ``count'' means using only the total number of detected events and ``pattern'' means the optimal test using the full space-time pattern. Hence the ratio of the numbers of rounds needed to reach the same error satisfies
$T_{\rm count}/T_{\rm pattern}\to D_{\rm pattern}/D_{\rm count}\ (\delta\to0)$, bounded above by
$D_{\rm event}/D_{\rm count}$.\hfill The chain of inequalities is \st{PROVED} (data-processing inequality for relative entropy/R\'enyi divergences and the Poisson closed forms); the asymptotic form of the ratio of rounds is \st{PROVED} under the marking model (M) of \S\ref{sec:design-theorems} (Theorem~\ref{thm:stein-rounds}; the error held at a fixed level is the one under $H_1$)
\end{proposition}

\begin{problem}[O4: experimental-design theorem]\label{prob:O4}
Under a budget constraint choose the code distance $d$, the number of rounds, and the (adaptive) strategy so as to maximize $D_{\rm pattern}(d,\text{strategy})$ or the Chernoff exponent;
and characterize the dependence of $D_{\rm pattern}/D_{\rm count}$ on the detector channel $W$.\hfill\st{PLAN}
\end{problem}
\noindent The numerics of P2 estimate $D_{\rm count}$ and $D_{\rm pattern}$ separately (by Monte Carlo or analytically), report $T_{\rm count}/T_{\rm pattern}$, and compare with Proposition~\ref{prop:poisson}
(the lower bound $\ln(1/\delta)/r$ on the number of rounds for any method).

\subsection{Experimental design under the marking model (O4)}\label{sec:design-theorems}

This subsection proves the asymptotic part of Proposition~\ref{prop:exponents}, gives an exact formula for the gain of pattern methods over counting, and shows that adaptive choice among experiments does not improve the Stein exponent. Time is measured in rounds.

\paragraph{Assumption (M): marked Poisson observation.} Under hypothesis $h\in\{0,1\}$, (i) events of class $a$ occur as independent Poisson processes of rate $r^{(h)}_a$ per round; (ii) each event independently produces one observed \emph{mark} $z$ in a finite set $\mathcal Z$ with probability $W(z\mid a)$ ($\sum_zW(z\mid a)\le1$; the deficit is the probability that the event is not seen); (iii) spurious marks from the background form an independent Poisson process of intensity $b(z)$ per round, the same under both hypotheses; (iv) the marks of different events are observed separately and are assigned to the round in which the event occurs. Put
\[
\lambda_h(z)=b(z)+\sum_ar^{(h)}_aW(z\mid a),\qquad\Lambda_h=\sum_z\lambda_h(z),\qquad\pi_h=\lambda_h/\Lambda_h,
\]
and assume $\lambda_0(z)>0\Leftrightarrow\lambda_1(z)>0$ and $\Lambda_0,\Lambda_1>0$.
(M) idealizes the clustering of P2: (iv) fails when events overlap in space-time, which is rare when $r$ times the event volume is small.

\begin{lemma}[Divergences of the observed data]\label{lem:poisson-kl}
Under (M) the data of $T$ rounds consist of independent counts $N_z\sim{\rm Poisson}(T\lambda_h(z))$, and given the counts the rounds of the marks are i.i.d.\ uniform (independently of $h$). Hence the relative entropy of the data of $T$ rounds under $H_1$ relative to $H_0$ is $T D_{\rm pattern}$, and that of the total count $N=\sum_zN_z$ is $TD_{\rm count}$, where
\[
D_{\rm pattern}=\sum_z\Big[\lambda_1\ln\frac{\lambda_1}{\lambda_0}-\lambda_1+\lambda_0\Big](z),\qquad D_{\rm count}=\Lambda_1\ln\frac{\Lambda_1}{\Lambda_0}-\Lambda_1+\Lambda_0 .
\]
The same holds for the Chernoff information, with $C_{\rm pattern}=\max_s\sum_z[s\lambda_1+(1-s)\lambda_0-\lambda_1^s\lambda_0^{1-s}](z)$ and its count analogue. $D_{\rm event}$ and $C$ of Proposition~\ref{prop:exponents} are the case $W={\rm id}$, $b=0$.\hfill\st{PROVED}
\end{lemma}
\begin{proof}
By the marking and superposition theorems for Poisson processes, the marks of type $z$ form independent Poisson processes of rates $\lambda_h(z)$. The law of the data is the product over $z$ of the count laws times the conditional law of the rounds, which does not depend on $h$; relative entropy and Chernoff coefficients of products factorize, and $D({\rm Poisson}(\mu_1)\|{\rm Poisson}(\mu_0))=\mu_1\ln(\mu_1/\mu_0)-\mu_1+\mu_0$, $\sum_n{\rm Pois}_{\mu_1}(n)^s{\rm Pois}_{\mu_0}(n)^{1-s}=\exp(\mu_1^s\mu_0^{1-s}-s\mu_1-(1-s)\mu_0)$. $N$ is Poisson with mean $T\Lambda_h$.
\end{proof}

\begin{theorem}[What pattern information is worth]\label{thm:pattern-gain}
Under (M),
\[
D_{\rm pattern}=D_{\rm count}+\Lambda_1\,D(\pi_1\|\pi_0).
\]
Hence $D_{\rm pattern}\ge D_{\rm count}$, with equality iff $\pi_1=\pi_0$ (the marks carry no information: $\lambda_1(z)/\lambda_0(z)$ does not depend on $z$), and
\[
\frac{D_{\rm pattern}}{D_{\rm count}}=1+\frac{\Lambda_1D(\pi_1\|\pi_0)}{D_{\rm count}}.
\]
$D_{\rm pattern}$ does not increase when the marks are coarse-grained, $z\mapsto z'$ with a kernel $K(z'\mid z)$ (equivalently $W\mapsto KW$, $b\mapsto Kb$).
In the weak-signal limit $H_1$: one extra class $a^\ast$ of rate $\epsilon$ on top of $H_0$ (so $\lambda_1=\lambda_0+\epsilon w$, $w=W(\cdot\mid a^\ast)$, $\|w\|_1>0$),
\[
\lim_{\epsilon\to0}\frac{D_{\rm pattern}}{D_{\rm count}}=1+\chi^2\big(\hat w\,\big\|\,\pi_0\big),\qquad\hat w=w/\|w\|_1,
\]
so the gain of pattern methods for weak signals is exactly one plus the $\chi^2$-distance between the mark distribution of the sought events and that of all other marks.\hfill\st{PROVED}
\end{theorem}
\begin{proof}
Write $\lambda_h=\Lambda_h\pi_h$: $\sum_z\lambda_1\ln(\lambda_1/\lambda_0)=\Lambda_1\ln(\Lambda_1/\Lambda_0)+\Lambda_1\sum_z\pi_1\ln(\pi_1/\pi_0)$, and $\sum_z(\lambda_0-\lambda_1)=\Lambda_0-\Lambda_1$. Equality: $D(\pi_1\|\pi_0)=0$ iff $\pi_1=\pi_0$.
Coarse-graining: $D_{\rm pattern}=\sum_z\lambda_0\,f(\lambda_1/\lambda_0)$ with $f(t)=t\ln t-t+1$ convex, so the log-sum (Jensen) inequality gives monotonicity under any kernel, exactly as for $f$-divergences.
Weak signal: $D({\rm Pois}(\mu+\epsilon w)\|{\rm Pois}(\mu))=\tfrac{\epsilon^2w^2}{2\mu}+O(\epsilon^3)$ term by term, so $D_{\rm pattern}=\tfrac{\epsilon^2}2\sum_zw(z)^2/\lambda_0(z)+O(\epsilon^3)$ and $D_{\rm count}=\tfrac{\epsilon^2}2\|w\|_1^2/\Lambda_0+O(\epsilon^3)$; the ratio tends to $\Lambda_0\sum_z\hat w^2/\lambda_0=\sum_z\hat w^2/\pi_0=1+\chi^2(\hat w\|\pi_0)$.
\end{proof}

\begin{theorem}[Stein form of the ratio of rounds]\label{thm:stein-rounds}
Assume (M) and fix $a\in(0,1)$. Among tests based on $T$ rounds whose error under $H_1$ is at most $a$, let $\beta_T$ be the smallest error under $H_0$, and let $T^\ast(\delta)=\min\{T\in\mathbb N:\beta_T\le\delta\}$. If $0<D<\infty$, then
\[
\lim_{\delta\to0}\frac{T^\ast(\delta)}{\ln(1/\delta)}=\frac1D,
\]
with $D=D_{\rm pattern}$ for tests that use all data and $D=D_{\rm count}$ for tests that use the total count only. Hence $T^\ast_{\rm count}(\delta)/T^\ast_{\rm pattern}(\delta)\to D_{\rm pattern}/D_{\rm count}$, as stated in Proposition~\ref{prop:exponents}.
For the Bayes error with fixed priors the same holds with the Chernoff informations $C_{\rm pattern},C_{\rm count}$; holding the error under $H_0$ fixed instead gives the reversed divergences $D(P_0\|P_1)$, for which Theorem~\ref{thm:pattern-gain} holds with the roles of $0$ and $1$ exchanged.\hfill\st{PROVED}
\end{theorem}
\begin{proof}
By Lemma~\ref{lem:poisson-kl} the data of $T$ rounds are $T$ i.i.d.\ blocks with per-block divergence $D$. Stein's lemma \cite[Thm.~11.8.3]{CT} (for every fixed $a\in(0,1)$, with the roles of the hypotheses as stated) gives $-\tfrac1T\ln\beta_T\to D$.
$\beta_T>0$ for every $T$, since the two laws are mutually absolutely continuous, and $\beta_T$ is nonincreasing (a test may ignore rounds); hence $T^\ast(\delta)\to\infty$ as $\delta\to0$.
Given $\eta>0$ choose $T_\eta$ with $e^{-(D+\eta)T}\le\beta_T\le e^{-(D-\eta)T}$ for $T\ge T_\eta$. Then $T^\ast(\delta)\le\max\{T_\eta,\lceil\ln(1/\delta)/(D-\eta)\rceil\}$, and for $\delta$ so small that $T^\ast(\delta)\ge T_\eta$, $\delta\ge\beta_{T^\ast}\ge e^{-(D+\eta)T^\ast}$, i.e.\ $T^\ast\ge\ln(1/\delta)/(D+\eta)$. Let $\eta\to0$.
The count version is the same argument applied to the Poisson count process. The Bayes version uses the Chernoff theorem \cite[Thm.~11.9.1]{CT} in place of Stein's lemma.
\end{proof}

\begin{theorem}[Adaptive designs do not beat the best single experiment (Stein)]\label{thm:adaptive}
Let $\mathcal I$ be a finite set of experiments (for example: run a patch of distance $d$ for one round with a given readout and clustering scheme). Experiment $i$ costs $c_i>0$ and yields data with law $P_{h,i}$ under $H_h$, with $D_i:=D(P_{1,i}\|P_{0,i})<\infty$. An adaptive design chooses the next experiment as a function of all past data (and private randomness), stops when the total cost would exceed a budget $\mathcal B$, and then decides. Let $D^\ast=\max_iD_i/c_i$.
If the design has error at most $a$ under $H_1$, its error $\alpha_0$ under $H_0$ satisfies
\[
\ln\frac1{\alpha_0}\;\le\;\frac{\mathcal B\,D^\ast+\ln2}{1-a}.
\]
Repeating the single experiment $i^\ast=\arg\max_iD_i/c_i$ achieves $\liminf_{\mathcal B\to\infty}\tfrac1{\mathcal B}\ln\tfrac1{\alpha_0}\ge D^\ast$ at every fixed $a\in(0,1)$. Hence the optimal exponent per unit cost, in the limit $a\to0$, is $D^\ast$, and it is attained by a non-adaptive design.\hfill\st{PROVED}
\end{theorem}
\begin{proof}
Let $Q_h$ be the law of the whole transcript (choices and data) under $H_h$. The number of steps is at most $\mathcal B/\min_ic_i$. By the chain rule for relative entropy, and because the choice at step $t$ has the same conditional law given the past under both hypotheses,
$D(Q_1\|Q_0)=\E_{Q_1}\sum_tD_{i_t}\le D^\ast\,\E_{Q_1}\sum_tc_{i_t}\le D^\ast\mathcal B$.
Let $x=Q_1[\text{decide }1]\ge1-a$ and $y=Q_0[\text{decide }1]=\alpha_0$. Data processing to the binary decision gives $D(Q_1\|Q_0)\ge x\ln\tfrac xy+(1-x)\ln\tfrac{1-x}{1-y}\ge x\ln\tfrac1y-\ln2$, which is the bound. Achievability is Stein's lemma for $\lfloor\mathcal B/c_{i^\ast}\rfloor$ i.i.d.\ uses of $i^\ast$.
\end{proof}

\noindent\emph{Reading for P2.} Theorems~\ref{thm:pattern-gain}--\ref{thm:adaptive} settle the two-hypothesis part of O4 (Problem~\ref{prob:O4}) under (M): the design question reduces to computing $D_{\rm pattern}(i)/c_i$ for each candidate $i$ (distance, readout, clustering) and taking the largest; the gain over counting is $\Lambda_1D(\pi_1\|\pi_0)$, which for weak signals is governed by $\chi^2(\hat w\|\pi_0)$; and the Poisson bound (Proposition~\ref{prop:poisson}) remains the floor for every design.
What is \emph{not} settled: (a) composite hypotheses (unknown shape, unknown rates), where adaptive designs can help and the relevant exponent is a minimax over the composite alternative; (b) the strong converse at fixed $a$ for adaptive designs (Theorem~\ref{thm:adaptive} uses the weak converse; the strong converse for classical adaptive discrimination is known in the literature but not reproduced here); (c) the Chernoff (Bayes) exponent for adaptive designs; (d) the computation of $D_{\rm pattern}(d)$ for the surface code, which needs the mark channel $W$ of P2; (e) the regime where (iv) fails (overlapping events).

\section{理论--数值对照表与开放问题}\label{sec:map}

原则:\textbf{每条定理都配一个数值证伪测试}。数值通过不证明定理,但违反则说明定理(或实现)有错;
数值结果也反过来决定下一个要证的命题。

\subsection{对照表}

{\footnotesize
\begin{longtable}{p{3.7cm} p{4.5cm} p{3.2cm} p{2.6cm}}
\toprule
理论陈述 & 数值检验 & 位置 / 里程碑 & 现状 \\
\midrule
\endfirsthead
\toprule
理论陈述 & 数值检验 & 位置 / 里程碑 & 现状 \\
\midrule
\endhead
定理~\ref{thm:monotone}、命题~\ref{prop:free-ops}\ (单调性) &
$\err(p,e)$ 对 $p,e$ 不降;叠加独立层不降 (C1,C2) &
\texttt{theory/checks}\ (精确 ML,3 个码) & 通过 \\
定理~\ref{thm:partial}\ (部分丢弃标记) &
$\err(p,e_0{+}\delta)\le\err(p'',e_0)$,约 $300$ 组随机参数/码 (C3) &
同上 & 通过 \\
定理~\ref{thm:preorder}\ (预序必要性) &
可达对均满足 $\err$ 同序 (C4);同时统计“不可达却同序”的比例(操作序严格大于自由序的有限 $d$ 证据) &
同上 & 通过;不可达对中约 15\% 仍同序 \\
引理~\ref{lem:union}、推论~\ref{cor:product}\ (联合上界) &
$\err\le\tfrac12\widetilde W(B)$ (C5) & 同上 & 通过 \\
命题~\ref{prop:genie}\ (先知下界) &
$\err_{d=3}\ge\varepsilon_{\rm genie}$ (C6);$d\ge5$ 用 TN-ML 复核 &
同上 / M1 & $d=3$ 通过 \\
定理~\ref{thm:rate}\ (汇率上界) &
有限 $d$:$R^{(d)}\le c$ (C7);\textbf{P1:$R_{e\to p}\le c(p_0,e_0)$,并报告 $\Delta=c-R$、$R_B$、$R_{\rm hash}$} &
\texttt{theory/checks};P1 / M2 & 有限 $d$ 通过;P1 待做 \\
观察~\ref{obs:sharp}\ ($[[5,1,3]]$ 取等) &
扫描 $e_0{=}0$ 处的 $R^{(d)}/c$:随机稳定子码库 &
\texttt{theory/checks}\ 扩展 & 3 个码已观察 \\
猜想~\ref{conj:universal}\ (码通用序 = 自由序) &
对不可达对 $(p,e)\to(p',e')$,在随机稳定子码库中搜索使 $\err$ \emph{反序}的码 &
\texttt{theory/checks}\ 扩展 & 未开始 \\
猜想~\ref{conj:HB}\ ($H_B$) &
P1 的 $R_\alpha$ 与解析的 $R_B$ 比较 &
P1 / M2 & 待 P1 \\
定理~\ref{thm:cc-cert}\ (码容量认证) &
(i) 模拟标定数据,检验 $\bar e,\bar p$ 的覆盖率 $\ge1-\delta$;(ii) 与 TN-ML 的 $\err_d$ 比较(证书必须 $\ge$) &
\texttt{experiments/}\ / M1 & 待做 \\
引理~\ref{lem:floor}\ (芯片事件地板) &
stim:$\pL\ge q_c(1-2^{-q})$ 且 $\propto r_cT$;地板与 $d$ 无关 &
P2 / M3--M4 & 待做 \\
命题~\ref{prop:poisson}\ (泊松下限) &
任何检验所需轮数 $\ge\ln((1-a)/\delta)/r$ &
P2 / M4 & 待做 \\
命题~\ref{prop:exponents}\ (指数链) &
估计 $D_{\rm count}$ 与 $D_{\rm pattern}$,验证 $T_{\rm count}/T_{\rm pattern}\to D_{\rm pattern}/D_{\rm count}$ &
P2 / M4 & 待做 \\
\bottomrule
\end{longtable}
}

\paragraph{精确小码的第二个用途。}
\texttt{exact\_small\_codes.py} 给出 $n\le9$ 的码上\emph{精确}的最大似然 $\err(p,e)$,可直接作为 M1(张量网络 ML 解码器)的 oracle:
$d=3$ 旋转表面码上 TN-ML 的结果必须与它在机器精度内一致(这是 README 中“合理性检验”之前更强的一道门)。

\paragraph{测试的分工:MWPM 不能证伪针对 ML 的上界。}
引理~\ref{lem:union}、定理~\ref{thm:cc-cert} 等是对 ML 的上界,而 MWPM 的失败率 $\ge\err$,
所以 MWPM 的数据既不能证实也不能证伪这些上界;必须用精确 ML(小 $d$)或 TN-ML。
而下界(地板、泊松下限、先知下界)对任何解码器成立,可以用 stim + MWPM 检验。

\subsection{开放问题清单}

\begin{enumerate}[leftmargin=2em]
\item \textbf{缺口问题}(问题~\ref{prob:gap}):$c-R_\alpha$ 来自码相关性还是自由操作不完备?
  \emph{先决数值}:P1 的 $R_\alpha$;\emph{理论路线}:猜想~\ref{conj:universal} 的 Blackwell/R\'enyi 路线。
\item \textbf{O1}(问题~\ref{prob:O1}):无标记突发下 $\err$ 对形状参数 $p_b,R,T_b$ 的单调性。
\item \textbf{O2}(问题~\ref{prob:O2}):突发/芯片事件污染下的参数置信上限。
\item \textbf{O3}(问题~\ref{prob:O3}):突发混合联合上界的显式估计(团簇展开)。
\item \textbf{O4}(问题~\ref{prob:O4}):实验设计定理。
\item \textbf{O5}:对\emph{次优解码器}(MWPM)的认证:需要针对该解码器的 Peierls 型上界,本笔记的上界只对 ML。
\item \textbf{O6}:观察~\ref{obs:sharp} 的证明:对哪些码类(完美码、传递码)单比特标记不改变 ML 判决?
\item \textbf{O7}:联合上界的收紧(用测得的小 $d$ 失败率校准),使认证在实际参数下不再空洞。
\end{enumerate}

\subsection{近期工作顺序(与 M0--M5 并行)}

\begin{center}
\begin{tabular}{lll}
\toprule
理论线 & 内容 & 与数值线的咬合 \\
\midrule
T0 & 框架、单调性、表示(本文 §\ref{sec:framework}--\ref{sec:repr})+ 精确小码检验 & M0/M1 的 oracle 与合理性检验 \\
T1 & 汇率上界、预序定理(§\ref{sec:exchange}) & M2:报告 $R\le c$、缺口 $\Delta$ \\
T2 & 猜想~\ref{conj:universal} 的反例搜索;观察~\ref{obs:sharp} 的证明 & 扩展 \texttt{checks};无需大计算 \\
T3 & 码容量认证 + O1--O3(§\ref{sec:cert}) & M3/M4:地板、参数识别 \\
T4 & 信息论界与实验设计(§\ref{sec:info}) & M4:$T_{\rm count}/T_{\rm pattern}$ 对照 \\
\bottomrule
\end{tabular}
\end{center}


\begin{thebibliography}{99}
\bibitem{DKLP} E.~Dennis, A.~Kitaev, A.~Landahl, J.~Preskill, \emph{Topological quantum memory}, J.\ Math.\ Phys.\ 43, 4452 (2002).
\bibitem{BSV} S.~Bravyi, M.~Suchara, A.~Vargo, \emph{Efficient algorithms for maximum likelihood decoding in the surface code}, Phys.\ Rev.\ A 90, 032326 (2014).
\bibitem{BV} S.~Bravyi, A.~Vargo, \emph{Simulation of rare events in quantum error correction}, Phys.\ Rev.\ A 88, 062308 (2013).
\bibitem{CF} C.~T.~Chubb, S.~T.~Flammia, \emph{Statistical mechanical models for quantum codes with correlated noise}, Ann.\ Inst.\ Henri Poincar\'e D 8, 269 (2021).
\bibitem{DZ} N.~Delfosse, G.~Z\'emor, \emph{Linear-time maximum likelihood decoding of surface codes over the quantum erasure channel}, Phys.\ Rev.\ Research 2, 033042 (2020).
\bibitem{Blackwell} D.~Blackwell, \emph{Equivalent comparisons of experiments}, Ann.\ Math.\ Stat.\ 24, 265 (1953).
\bibitem{MPST} X.~Mu, L.~Pomatto, P.~Strack, O.~Tamuz, \emph{From Blackwell dominance in large samples to R\'enyi divergences and back again}, Econometrica 89, 475 (2021).
\bibitem{CT} T.~Cover, J.~Thomas, \emph{Elements of Information Theory}, 2nd ed., Wiley (2006), Ch.~11.
\bibitem{LeCam} L.~Le Cam, \emph{Asymptotic Methods in Statistical Decision Theory}, Springer (1986).
\end{thebibliography}
\end{document}
